\documentclass[12pt,a4paper]{article}
\usepackage[utf8]{inputenc}
\usepackage[T1]{fontenc}
\usepackage[french]{babel}
\usepackage{geometry}
\usepackage{graphicx}
\usepackage{enumitem}
\usepackage{longtable}
\usepackage{array}
\usepackage{xcolor}
\usepackage{titlesec}
\usepackage{fancyhdr}
\usepackage{hyperref}
\usepackage{colortbl}
\usepackage{booktabs}
\usepackage{amsmath}
\usepackage{amssymb}
\usepackage{float}
\usepackage{mathpazo} % Palatino font for a premium look
\usepackage{tcolorbox} % For beautiful explanation boxes
\tcbuselibrary{skins,breakable}

\geometry{margin=2.5cm}
\setlength{\headheight}{15pt}

\definecolor{bleu}{HTML}{2B6CB0}
\definecolor{vertclair}{HTML}{48BB78}
\definecolor{grisclair}{HTML}{F7FAFC}

\titleformat{\section}
  {\Large\bfseries\color{bleu}}{\thesection.}{0.4em}{}

\titleformat{\subsection}
  {\large\bfseries\color{bleu}}{\thesubsection}{0.4em}{}

\titleformat{\subsubsection}
  {\normalsize\bfseries\color{bleu}}{\thesubsubsection}{0.4em}{}

\pagestyle{fancy}
\fancyhf{}
\fancyhead[L]{\small Chapitre Conception -- Nexora}
\fancyhead[R]{\small EMSI 2025-2026}
\fancyfoot[C]{\thepage}
\renewcommand{\headrulewidth}{0.4pt}

\hypersetup{colorlinks=true, linkcolor=bleu, urlcolor=bleu}

\begin{document}

% ======================================================================
\begin{titlepage}
\centering
\vspace*{1.5cm}

\includegraphics[width=0.45\textwidth]{emsilogo.jpeg}\\[2.5cm]

\rule{\linewidth}{0.6pt}\\[0.6cm]
{\LARGE\textbf{Chapitre -- Conception}}\\[0.4cm]
{\Huge\bfseries\color{bleu} Conception Fonctionnelle}\\[0.4cm]
{\Large\textit{Modélisation Métier, Acteurs et Parcours Utilisateurs}}\\[0.4cm]
\rule{\linewidth}{0.6pt}

\vspace{2.5cm}

\begin{tabular}{cc}
    \multicolumn{2}{c}{\large\textbf{Réalisé par :}} \\[0.3cm]
    Soubaibi Omar \& Yassine Baizi \\[1cm]
    \multicolumn{2}{c}{\large\textbf{Encadré par :}} \\[0.3cm]
    \multicolumn{2}{c}{Mme. OUNACER SOUMAYA} \\
\end{tabular}

\vfill
\textbf{Filière :} 4IADM G5-A \\[0.2cm]
\textbf{Année Universitaire :} 2025--2026
\end{titlepage}

\tableofcontents
\newpage

% ======================================================================
\section{Introduction du chapitre}
% ======================================================================

La phase de conception constitue une étape charnière dans le cycle de vie de tout projet de développement logiciel. Faisant suite à l'analyse approfondie des besoins et à la définition des exigences fonctionnelles et non fonctionnelles, cette étape vise à traduire les attentes exprimées en une architecture logique et conceptuelle rigoureuse. Le présent chapitre est ainsi dédié à la modélisation fonctionnelle et structurelle de la plateforme Nexora, une solution innovante destinée à la cartographie des connaissances et à la découverte d'experts au sein des organisations complexes.

L'approche de conception adoptée s'appuie sur les principes de la modélisation orientée objet, permettant d'abstraire la complexité inhérente au traitement des relations humaines et professionnelles. Ce chapitre s'articulera autour de plusieurs axes fondamentaux. Dans un premier temps, nous procéderons à l'identification exhaustive des acteurs interagissant avec le système. Ensuite, nous détaillerons les cas d'utilisation, véritables reflets des interactions fonctionnelles, avant de proposer une description structurée des scénarios les plus critiques. Par la suite, une modélisation statique sera présentée à travers le diagramme de classes métier, illustrant les entités centrales du domaine et leurs interrelations. Enfin, la projection visuelle de cette conception sera matérialisée par la description des maquettes fonctionnelles, offrant un aperçu tangible de l'expérience utilisateur finale, avant de conclure sur les perspectives de réalisation.

% ======================================================================
\section{Identification des acteurs}
% ======================================================================

Dans le cadre de la modélisation d'un système interactif, un acteur représente toute entité externe (utilisateur humain ou système tiers) amenée à interagir directement avec la plateforme pour accomplir une ou plusieurs tâches spécifiques. Pour le système Nexora, nous avons identifié plusieurs profils d'acteurs, chacun doté de prérogatives et d'objectifs distincts, reflétant la diversité des parties prenantes au sein d'une organisation moderne.

\begin{longtable}{|p{3cm}|p{10.5cm}|}
\toprule
\rowcolor{bleu!15}
\textbf{Acteur} & \textbf{Description et Responsabilités} \\
\midrule
\endfirsthead
\toprule
\rowcolor{bleu!15}
\textbf{Acteur} & \textbf{Description et Responsabilités} \\
\midrule
\endhead
\textbf{Employé} & Il est le premier acteur et le plus central. Il s'agit de l'utilisateur final qui navigue sur la plateforme pour valoriser son profil, chercher de l'information ou identifier des collaborateurs qualifiés. L'Employé est à la fois consommateur et producteur de connaissances. Son objectif principal est de se repérer dans le tissu intellectuel de l'entreprise et de développer son parcours professionnel continu. \\
\midrule
\textbf{Manager RH} & Ce deuxième acteur clé dispose d'une vision plus stratégique et agrégée. Il interagit avec le système non pas pour une recherche ponctuelle d'information, mais pour analyser des tendances, anticiper les besoins en formation et orchestrer la composition des équipes. Le Manager RH utilise la plateforme comme un outil d'aide à la décision stratégique pour la gestion prévisionnelle des emplois et des compétences. \\
\midrule
\textbf{Administrateur} & Ce troisième acteur humain possède un rôle strictement opérationnel et global. Il est le garant du bon fonctionnement, de la cohérence et de la sécurité du système. Ses interactions visent à paramétrer la plateforme, gérer les droits d'accès des différents utilisateurs, et s'assurer de la validité des données qui alimentent le système, notamment par la supervision des flux documentaires et des référentiels de l'organisation. \\
\midrule
\textbf{Veda (Assistant)} & Enfin, la modélisation conceptuelle de Nexora intègre un acteur secondaire non humain : le chatbot Veda. Conçu comme un agent conversationnel intelligent, Veda agit en tant qu'assistant virtuel proactif. Il se comporte conceptuellement comme un acteur externe vis-à-vis des autres logiques métiers, en fournissant des réponses dynamiques, en orientant les recherches des utilisateurs et en facilitant l'interaction en langage naturel. \\
\bottomrule
\end{longtable}

% ======================================================================
\section{Diagramme des cas d'utilisation}
% ======================================================================

Le diagramme des cas d'utilisation offre une vue d'ensemble de la cartographie fonctionnelle du système, en illustrant les services rendus par Nexora à ses différents acteurs. Ce modèle permet de circonscrire précisément le périmètre d'action de chaque profil utilisateur au sein de la plateforme.

\begin{figure}[H]
    \centering
    \includegraphics[width=0.75\textwidth]{plantuml/output/01_use_case.png}
    \caption{Diagramme de Cas d'Utilisation -- Nexora}
    \label{fig:usecase}
\end{figure}

Pour \textbf{l'Employé}, les interactions sont principalement orientées vers la consultation, l'exploration et l'apprentissage. 

\begin{longtable}{|p{1cm}|p{4.5cm}|p{8cm}|}
\toprule
\rowcolor{bleu!15}
\textbf{ID} & \textbf{Cas d'utilisation} & \textbf{Description et Flux} \\
\midrule
\endfirsthead
\toprule
\rowcolor{bleu!15}
\textbf{ID} & \textbf{Cas d'utilisation} & \textbf{Description et Flux} \\
\midrule
\endhead
UC01 & Consulter son profil & Consultation de son propre profil de compétences pour visualiser son évolution et les validations obtenues par ses pairs. \\
\midrule
UC02 & Explorer la carte & Exploration de la carte interactive des connaissances, offrant une vue spatiale et relationnelle du savoir réparti. \\
\midrule
UC03 & Rechercher un expert & Déclenchée lorsqu'un besoin de résolution de problème complexe ou de mentorat se présente. \\
\midrule
UC04 & Recevoir des recos. & Le système fournit proactivement des recommandations de compétences pour guider son développement personnel. \\
\midrule
UC05 & Interagir avec Veda & Interaction de manière conversationnelle avec l'assistant Veda pour des requêtes plus directes ou intuitives. \\
\bottomrule
\end{longtable}

Le \textbf{Manager RH}, quant à lui, dispose de cas d'utilisation résolument orientés vers l'analyse globale et la gestion stratégique.

\begin{longtable}{|p{1cm}|p{4.5cm}|p{8cm}|}
\toprule
\rowcolor{bleu!15}
\textbf{ID} & \textbf{Cas d'utilisation} & \textbf{Description et Flux} \\
\midrule
\endfirsthead
\toprule
\rowcolor{bleu!15}
\textbf{ID} & \textbf{Cas d'utilisation} & \textbf{Description et Flux} \\
\midrule
\endhead
UC06 & Visualiser les lacunes & Visualisation des lacunes en compétences à l'échelle d'un département ou de l'entreprise entière. \\
\midrule
UC07 & Identifier les experts & Identification des experts clés pour les valoriser, organiser des transferts ou les affecter à des missions stratégiques. \\
\midrule
UC08 & Tendances émergentes & Consultation des tendances de compétences émergentes pour adapter la politique de recrutement et de formation. \\
\midrule
UC09 & Gérer les équipes & Gestion de la composition des équipes en s'appuyant sur les recommandations algorithmiques du système. \\
\bottomrule
\end{longtable}

L'\textbf{Administrateur} possède des privilèges transverses, focalisés sur la maintenance, la configuration et la modération.

\begin{longtable}{|p{1cm}|p{4.5cm}|p{8cm}|}
\toprule
\rowcolor{bleu!15}
\textbf{ID} & \textbf{Cas d'utilisation} & \textbf{Description et Flux} \\
\midrule
\endfirsthead
\toprule
\rowcolor{bleu!15}
\textbf{ID} & \textbf{Cas d'utilisation} & \textbf{Description et Flux} \\
\midrule
\endhead
UC10 & Gérer les utilisateurs & Gestion complète du cycle de vie des utilisateurs, allant de la création des accès à la modification des rôles. \\
\midrule
UC11 & Classifier les documents & Supervision de la classification automatisée des documents intégrés à la plateforme, en corrigeant ou validant les catégories. \\
\midrule
UC12 & Superviser le système & Orchestration de la supervision globale de la plateforme, garantissant son intégrité et la conformité des données. \\
\bottomrule
\end{longtable}

% ======================================================================
\section{Description des cas d'utilisation principaux}
% ======================================================================

Afin de détailler la dynamique interne du système et de comprendre la chronologie des interactions, il convient de spécifier textuellement le déroulement logique des fonctionnalités les plus névralgiques de la plateforme.

\begin{tcolorbox}[colback=bleu!5,colframe=bleu,title=Cas d'utilisation 1 : Recherche d'expert]
\textbf{Acteur principal :} Employé. \\
\textbf{Préconditions :} L'Employé doit être authentifié avec succès et disposer d'un profil actif sur la plateforme. \\
\textbf{Scénario nominal :}
\begin{enumerate}
    \item L'Employé accède à l'interface de recherche globale et saisit une thématique, un domaine de compétence pointu ou la description textuelle d'une problématique spécifique.
    \item Le système réceptionne la requête et lance une analyse sémantique approfondie.
    \item Le système parcourt le réseau de relations entre les collaborateurs, examine les projets passés et scrute les documents publiés pour identifier les profils les plus pertinents.
    \item Le système évalue et pondère le niveau d'expertise de chaque profil trouvé selon sa position dans le réseau de connaissances de l'entreprise.
    \item Le système retourne à l'Employé une liste ordonnée d'experts potentiels, en mettant en exergue les raisons précises de leur sélection (par exemple, la participation à un projet très similaire ou la possession d'une compétence certifiée de haut niveau).
    \item L'Employé sélectionne un profil pour consulter ses détails approfondis et initie, le cas échéant, une demande de prise de contact formelle.
\end{enumerate}
\textbf{Postconditions :} L'Employé a obtenu une sélection d'experts qualifiés répondant à son besoin, et une potentielle mise en relation professionnelle a été initiée et tracée par le système.
\end{tcolorbox}

\begin{tcolorbox}[colback=bleu!5,colframe=bleu,title=Cas d'utilisation 2 : Visualisation de la carte des connaissances]
\textbf{Acteur principal :} Employé (ou Manager RH). \\
\textbf{Préconditions :} L'utilisateur est authentifié et possède les droits de consultation des cartographies. \\
\textbf{Scénario nominal :}
\begin{enumerate}
    \item L'utilisateur navigue vers le module dédié à la cartographie organisationnelle.
    \item Le système compile les données relationnelles à jour et charge une représentation visuelle, dynamique et tridimensionnelle des savoirs de l'organisation.
    \item L'utilisateur observe un écosystème de nœuds interactifs représentant des individus, des domaines d'expertise, des projets et des départements.
    \item L'utilisateur manipule la carte (zoom, rotation, translation) pour explorer différentes zones de l'entreprise.
    \item Lorsqu'il sélectionne un nœud spécifique, le système met instantanément en surbrillance ses connexions directes et indirectes, révélant ainsi les réseaux invisibles (qui connaît qui, qui détient quelle expertise rare).
    \item L'utilisateur applique des filtres thématiques ou structurels pour restreindre l'affichage à une entité précise de l'entreprise afin de focaliser son analyse.
\end{enumerate}
\textbf{Postconditions :} L'utilisateur a acquis une compréhension spatiale, relationnelle et intuitive de la distribution des expertises au sein de son environnement de travail.
\end{tcolorbox}

\begin{tcolorbox}[colback=bleu!5,colframe=bleu,title=Cas d'utilisation 3 : Détection des lacunes en compétences]
\textbf{Acteur principal :} Manager RH. \\
\textbf{Préconditions :} Le Manager RH est authentifié avec un profil disposant de privilèges d'accès aux données analytiques globales et stratégiques. \\
\textbf{Scénario nominal :}
\begin{enumerate}
    \item Le Manager RH accède au tableau de bord de pilotage analytique et sélectionne l'outil de diagnostic des écarts de compétences.
    \item Il paramètre un périmètre d'étude précis, tel qu'une équipe de projet, un pôle d'activité entier ou une filiale.
    \item Il définit ensuite une cible stratégique ou un référentiel requis (par exemple, les compétences nécessaires pour lancer une nouvelle ligne de produits).
    \item Le système croise et compare l'inventaire actuel des aptitudes du groupe sélectionné avec le référentiel cible projeté.
    \item Le système identifie automatiquement les déficits, qu'ils soient d'ordre quantitatif (manque d'effectif) ou qualitatif (manque de séniorité), et génère une représentation graphique limpide de ces écarts.
    \item Le système conclut en proposant des plans d'action prédictifs, suggérant des parcours de formation interne croisés ou la nécessité d'une campagne de recrutement ciblée.
\end{enumerate}
\textbf{Postconditions :} Un rapport décisionnel détaillé cartographiant les vulnérabilités en capital humain est généré, accompagné de recommandations algorithmiques visant à orienter la stratégie de développement des talents.
\end{tcolorbox}

% ======================================================================
\section{Diagramme de classes métier}
% ======================================================================

La modélisation de la structure de l'information est une étape fondamentale pour asseoir la robustesse à long terme du système. Le diagramme de classes métier offre une cartographie conceptuelle des entités manipulées par Nexora. Ce modèle est volontairement expurgé de toute considération technique ou d'implémentation afin de se concentrer exclusivement sur la sémantique et la taxonomie de l'organisation.

\begin{figure}[H]
    \centering
    \includegraphics[width=0.85\textwidth]{plantuml/output/02_class_diagram.png}
    \caption{Diagramme de Classes -- Nexora}
    \label{fig:classes}
\end{figure}

\begin{longtable}{|p{2.5cm}|p{11cm}|}
\toprule
\rowcolor{bleu!15}
\textbf{Classe} & \textbf{Rôle et Informations Clés} \\
\midrule
\endfirsthead
\toprule
\rowcolor{bleu!15}
\textbf{Classe} & \textbf{Rôle et Informations Clés} \\
\midrule
\endhead
\textbf{Utilisateur} & Entité centrale et névralgique de ce modèle. Cette classe généralise les différents acteurs humains du système et porte les informations fondamentales d'identité, de rôle et de rattachement. \\
\midrule
\textbf{Compétence} & L'association avec l'utilisateur se caractérise par un niveau de maîtrise avéré et un historique de validation. Une compétence métier est détenue et partagée par un ensemble d'utilisateurs. \\
\midrule
\textbf{Projet} & Entité temporelle rassemblant plusieurs utilisateurs autour d'un objectif de production commun. Cette association permet au système de tracer et de contextualiser l'expérience pratique réelle. \\
\midrule
\textbf{Document} & Incarne les livrables, rapports, méthodologies ou notes internes constituant le capital intellectuel explicite de l'entreprise, auquel les utilisateurs contribuent. \\
\midrule
\textbf{Organisation} & Entité administrative qui regroupe les utilisateurs en départements, définissant ainsi la hiérarchie formelle et les frontières structurelles. \\
\midrule
\textbf{Badge} & Intégré pour répondre aux enjeux de motivation et de ludification (gamification). Décerné en reconnaissance de l'implication continue ou de la validation croisée. \\
\bottomrule
\end{longtable}

\begin{tcolorbox}[colback=vertclair!10,colframe=vertclair,title=Dynamique des relations sémantiques]
L'entité Utilisateur maintient une relation étroite et structurante avec la classe Compétence. Parallèlement, l'entité Utilisateur s'inscrit dans une dynamique d'accomplissement à travers son association avec la classe Projet. Les entités Projet et Document sont, à leur tour, intimement liées à la classe Compétence : un projet exige la mobilisation d'un ensemble de compétences, tandis qu'un document est classifié sémantiquement comme traitant certains domaines d'expertise.
\end{tcolorbox}

% ======================================================================
\section{Diagrammes de séquence}
% ======================================================================

Les diagrammes de séquence modélisent la chronologie des échanges d'informations entre les acteurs et les différents composants du système lors de la réalisation des processus clés.

\subsection{Séquence : Authentification sécurisée}

\begin{figure}[H]
    \centering
    \includegraphics[width=0.85\textwidth]{plantuml/output/03_sequence_auth.png}
    \caption{Diagramme de Séquence -- Connexion Sécurisée}
    \label{fig:seq_auth}
\end{figure}

\begin{tcolorbox}[colback=bleu!5,colframe=bleu,title=Flux logique de l'authentification]
\begin{enumerate}[leftmargin=*]
    \item \textbf{Requête initiale :} L'utilisateur soumet ses identifiants via l'interface d'accès.
    \item \textbf{Vérification d'identité :} L'interface transmet la demande au gestionnaire d'authentification central.
    \item \textbf{Validation :} Le gestionnaire consulte le référentiel des profils pour vérifier la validité des informations fournies.
    \item \textbf{Attribution des droits :} Une fois l'identité confirmée, une session sécurisée est établie et un profil de droits (rôle) est attribué, garantissant que l'utilisateur n'accède qu'aux fonctionnalités correspondant à ses prérogatives.
\end{enumerate}
\end{tcolorbox}

\subsection{Séquence : Recommandation d'expert}

\begin{figure}[H]
    \centering
    \includegraphics[width=0.85\textwidth]{plantuml/output/04_sequence_recommendation.png}
    \caption{Diagramme de Séquence -- Recherche du Meilleur Profil}
    \label{fig:seq_recommend}
\end{figure}

\begin{tcolorbox}[colback=bleu!5,colframe=bleu,title=Flux logique de la recommandation]
\begin{enumerate}[leftmargin=*]
    \item \textbf{Formulation du besoin :} L'employé formule une requête de recherche d'expert en langage naturel.
    \item \textbf{Extraction des profils :} Le moteur de recherche interroge la base de connaissances pour extraire l'ensemble des profils potentiellement pertinents.
    \item \textbf{Analyse intelligente :} Le moteur d'analyse sémantique calcule le degré de similarité entre la requête de l'utilisateur et le profil complet de chaque expert (compétences, projets, publications).
    \item \textbf{Restitution :} Le système consolide les résultats, trie les profils par score de pertinence décroissant et renvoie la liste ordonnée à l'employé.
\end{enumerate}
\end{tcolorbox}

\subsection{Séquence : Évaluation de l'influence (Classement)}

\begin{figure}[H]
    \centering
    \includegraphics[width=0.85\textwidth]{plantuml/output/05_sequence_pagerank.png}
    \caption{Diagramme de Séquence -- Classement des Experts}
    \label{fig:seq_pagerank}
\end{figure}

\begin{tcolorbox}[colback=bleu!5,colframe=bleu,title=Flux logique du classement des experts]
\begin{enumerate}[leftmargin=*]
    \item \textbf{Déclenchement :} Le système lance une requête d'évaluation globale des profils.
    \item \textbf{Exploration du réseau :} Le moteur d'évaluation parcourt le graphe relationnel pour mesurer la centralité de chaque collaborateur (nombre et qualité de ses connexions au sein de l'entreprise).
    \item \textbf{Pondération :} Cette mesure d'influence par le réseau est ensuite pondérée et combinée avec des critères objectifs tels que l'ancienneté, la diversité des compétences ou le nombre de projets validés.
    \item \textbf{Classement final :} Le classement consolidé est généré et restitué, permettant de mettre en lumière les véritables leaders d'opinion souvent invisibles dans la structure hiérarchique classique.
\end{enumerate}
\end{tcolorbox}

% ======================================================================
\section{Maquettes fonctionnelles (wireframes)}
% ======================================================================

La transition de la modélisation abstraite vers une proposition ergonomique et centrée sur l'utilisateur se matérialise par la conception des maquettes fonctionnelles. Ces représentations permettent de figer la disposition spatiale de l'information, la hiérarchie visuelle et le parcours de navigation logique, garantissant ainsi une expérience utilisateur optimale, intuitive et exempte de frictions.

\subsection{Tableau de bord personnel de l'employé}
Il se présente comme la porte d'entrée quotidienne de l'utilisateur. Conçu de manière visuellement épurée pour éviter la surcharge cognitive, il met en évidence, dans sa partie supérieure, une synthèse valorisante du profil : résumé de l'expertise principale, niveau d'engagement global et derniers badges honorifiques acquis. Le cœur de l'interface est consacré à un fil d'actualité intelligent et personnalisé, suggérant des lectures internes pertinentes ou des recommandations de formations calculées sur la base des aspirations et des lacunes identifiées de l'employé. Un panneau latéral persistant permet un accès rapide aux requêtes de recherche récentes et regroupe les notifications de mise en relation, assurant ainsi une prise en main immédiate et un engagement continu.

\subsection{Interface de la carte interactive des connaissances}
Elle constitue l'innovation visuelle majeure et la signature de la plateforme. L'écran entier est alloué à une visualisation immersive et fluide où chaque élément de l'entreprise flotte dans un espace défini. Les utilisateurs, identifiés par des avatars circulaires photographiques, sont reliés par des segments graphiques aux concepts ou compétences qu'ils maîtrisent, représentés par des nœuds géométriques distincts. Des panneaux de contrôle, situés de manière minimaliste et non intrusive sur les bordures de l'écran, offrent la possibilité d'ajuster la granularité de la vue, d'isoler un département particulier ou d'isoler un nœud spécifique. Lors du survol d'un élément par le curseur, une infobulle élégante et réactive surgit instantanément, dévoilant les informations contextuelles essentielles sans nécessiter de quitter la vue panoramique.

\subsection{Page de recherche et de découverte d'experts}
Elle a été pensée pour maximiser l'efficacité de l'investigation. Elle s'articule autour d'une barre de saisie proéminente, positionnée au centre, invitant l'utilisateur à exprimer ses requêtes en langage naturel de façon libre. En périphérie, des filtres avancés permettent de raffiner méthodiquement la recherche par localisation géographique, département d'appartenance ou niveau d'ancienneté. Les résultats sont restitués sous forme de fiches profil synthétiques et aérées. Chaque fiche met en exergue non seulement le score d'adéquation de la personne avec la requête initiale, mais affiche également un "chemin relationnel" visuel illustrant concrètement comment l'utilisateur actuel est connecté à cet expert potentiel (par le biais d'un ancien manager commun ou d'une participation croisée à un projet), favorisant ainsi grandement le premier contact formel.

\subsection{Interface d'analyse des compétences pour le manager RH}
Elle adopte une approche résolument décisionnelle et macroscopique. Elle est structurée autour de tableaux de bord composés de graphiques de synthèse et de jauges de performance de l'organisation. Une section prépondérante affiche une matrice dynamique des écarts de compétences, mettant en parallèle la situation capacitrice actuelle avec les objectifs stratégiques à l'aide d'une cartographie colorimétrique intuitive. Des listes de données interactives répertorient les compétences émergentes, systématiquement accompagnées d'indicateurs de vélocité illustrant l'accélération de leur adoption au sein des équipes. Cette interface se veut avant tout sobre, hautement dense en informations stratégiques utiles, et propose des appels à l'action clairs pour exporter des rapports formels ou déclencher des processus organisationnels de planification des ressources.

% ======================================================================
\section{Conclusion du chapitre}
% ======================================================================

En définitive, cette phase de conception fonctionnelle et structurelle de la plateforme Nexora a permis de poser avec rigueur les fondements logiques indispensables à la réalisation d'un outil pérenne, scalable et profondément centré sur l'utilisateur. En identifiant de manière exhaustive les acteurs et en modélisant minutieusement les cas d'utilisation, nous avons pu garantir que chaque fonctionnalité envisagée répond à un besoin métier réel, qualifié et documenté de l'entreprise, qu'il s'agisse de la résolution d'un problème ponctuel par la recherche d'experts ou de la prise de hauteur par l'analyse stratégique des compétences à l'échelle organisationnelle. 

Le diagramme de classes métier a, pour sa part, défini une architecture sémantique de l'information solide. Cette modélisation s'est avérée capable de retranscrire avec une grande finesse la complexité des interrelations professionnelles et du savoir, sans se heurter, à ce stade du cycle de vie, aux contraintes de faisabilité technique. Enfin, la projection au travers des maquettes fonctionnelles a concrétisé cette vision purement abstraite en proposant une interface de navigation fluide et engageante, dimensionnée pour favoriser une adoption rapide par l'ensemble des collaborateurs de l'organisation. 

Ces choix de conception, tant sur le plan de la logique conceptuelle que de l'ergonomie, constituent aujourd'hui un socle stable, unifié et cohérent. Ils balisent ainsi naturellement et méthodiquement la voie vers la phase subséquente du projet : le chapitre de réalisation, qui s'attachera dès lors à traduire ces modèles conceptuels abstraits en une architecture logicielle concrète, déployable et opérationnelle.

\end{document}
